% Options for packages loaded elsewhere
\PassOptionsToPackage{unicode}{hyperref}
\PassOptionsToPackage{hyphens}{url}
\documentclass[
]{article}
\usepackage{xcolor}
\usepackage{amsmath,amssymb}
\setcounter{secnumdepth}{-\maxdimen} % remove section numbering
\usepackage{iftex}
\ifPDFTeX
  \usepackage[T1]{fontenc}
  \usepackage[utf8]{inputenc}
  \usepackage{textcomp} % provide euro and other symbols
\else % if luatex or xetex
  \usepackage{unicode-math} % this also loads fontspec
  \defaultfontfeatures{Scale=MatchLowercase}
  \defaultfontfeatures[\rmfamily]{Ligatures=TeX,Scale=1}
\fi
\usepackage{lmodern}
\ifPDFTeX\else
  % xetex/luatex font selection
\fi
% Use upquote if available, for straight quotes in verbatim environments
\IfFileExists{upquote.sty}{\usepackage{upquote}}{}
\IfFileExists{microtype.sty}{% use microtype if available
  \usepackage[]{microtype}
  \UseMicrotypeSet[protrusion]{basicmath} % disable protrusion for tt fonts
}{}
\makeatletter
\@ifundefined{KOMAClassName}{% if non-KOMA class
  \IfFileExists{parskip.sty}{%
    \usepackage{parskip}
  }{% else
    \setlength{\parindent}{0pt}
    \setlength{\parskip}{6pt plus 2pt minus 1pt}}
}{% if KOMA class
  \KOMAoptions{parskip=half}}
\makeatother
\usepackage{color}
\usepackage{fancyvrb}
\newcommand{\VerbBar}{|}
\newcommand{\VERB}{\Verb[commandchars=\\\{\}]}
\DefineVerbatimEnvironment{Highlighting}{Verbatim}{commandchars=\\\{\}}
% Add ',fontsize=\small' for more characters per line
\newenvironment{Shaded}{}{}
\newcommand{\AlertTok}[1]{\textcolor[rgb]{1.00,0.00,0.00}{\textbf{#1}}}
\newcommand{\AnnotationTok}[1]{\textcolor[rgb]{0.38,0.63,0.69}{\textbf{\textit{#1}}}}
\newcommand{\AttributeTok}[1]{\textcolor[rgb]{0.49,0.56,0.16}{#1}}
\newcommand{\BaseNTok}[1]{\textcolor[rgb]{0.25,0.63,0.44}{#1}}
\newcommand{\BuiltInTok}[1]{\textcolor[rgb]{0.00,0.50,0.00}{#1}}
\newcommand{\CharTok}[1]{\textcolor[rgb]{0.25,0.44,0.63}{#1}}
\newcommand{\CommentTok}[1]{\textcolor[rgb]{0.38,0.63,0.69}{\textit{#1}}}
\newcommand{\CommentVarTok}[1]{\textcolor[rgb]{0.38,0.63,0.69}{\textbf{\textit{#1}}}}
\newcommand{\ConstantTok}[1]{\textcolor[rgb]{0.53,0.00,0.00}{#1}}
\newcommand{\ControlFlowTok}[1]{\textcolor[rgb]{0.00,0.44,0.13}{\textbf{#1}}}
\newcommand{\DataTypeTok}[1]{\textcolor[rgb]{0.56,0.13,0.00}{#1}}
\newcommand{\DecValTok}[1]{\textcolor[rgb]{0.25,0.63,0.44}{#1}}
\newcommand{\DocumentationTok}[1]{\textcolor[rgb]{0.73,0.13,0.13}{\textit{#1}}}
\newcommand{\ErrorTok}[1]{\textcolor[rgb]{1.00,0.00,0.00}{\textbf{#1}}}
\newcommand{\ExtensionTok}[1]{#1}
\newcommand{\FloatTok}[1]{\textcolor[rgb]{0.25,0.63,0.44}{#1}}
\newcommand{\FunctionTok}[1]{\textcolor[rgb]{0.02,0.16,0.49}{#1}}
\newcommand{\ImportTok}[1]{\textcolor[rgb]{0.00,0.50,0.00}{\textbf{#1}}}
\newcommand{\InformationTok}[1]{\textcolor[rgb]{0.38,0.63,0.69}{\textbf{\textit{#1}}}}
\newcommand{\KeywordTok}[1]{\textcolor[rgb]{0.00,0.44,0.13}{\textbf{#1}}}
\newcommand{\NormalTok}[1]{#1}
\newcommand{\OperatorTok}[1]{\textcolor[rgb]{0.40,0.40,0.40}{#1}}
\newcommand{\OtherTok}[1]{\textcolor[rgb]{0.00,0.44,0.13}{#1}}
\newcommand{\PreprocessorTok}[1]{\textcolor[rgb]{0.74,0.48,0.00}{#1}}
\newcommand{\RegionMarkerTok}[1]{#1}
\newcommand{\SpecialCharTok}[1]{\textcolor[rgb]{0.25,0.44,0.63}{#1}}
\newcommand{\SpecialStringTok}[1]{\textcolor[rgb]{0.73,0.40,0.53}{#1}}
\newcommand{\StringTok}[1]{\textcolor[rgb]{0.25,0.44,0.63}{#1}}
\newcommand{\VariableTok}[1]{\textcolor[rgb]{0.10,0.09,0.49}{#1}}
\newcommand{\VerbatimStringTok}[1]{\textcolor[rgb]{0.25,0.44,0.63}{#1}}
\newcommand{\WarningTok}[1]{\textcolor[rgb]{0.38,0.63,0.69}{\textbf{\textit{#1}}}}
\usepackage{longtable,booktabs,array}
\newcounter{none} % for unnumbered tables
\usepackage{calc} % for calculating minipage widths
% Correct order of tables after \paragraph or \subparagraph
\usepackage{etoolbox}
\makeatletter
\patchcmd\longtable{\par}{\if@noskipsec\mbox{}\fi\par}{}{}
\makeatother
% Allow footnotes in longtable head/foot
\IfFileExists{footnotehyper.sty}{\usepackage{footnotehyper}}{\usepackage{footnote}}
\makesavenoteenv{longtable}
\setlength{\emergencystretch}{3em} % prevent overfull lines
\providecommand{\tightlist}{%
  \setlength{\itemsep}{0pt}\setlength{\parskip}{0pt}}
\usepackage{bookmark}
\IfFileExists{xurl.sty}{\usepackage{xurl}}{} % add URL line breaks if available
\urlstyle{same}
\hypersetup{
  pdftitle={ECL: A Deterministic Cross-Runtime Execution IR with Replayable Semantics and Dependency Interface},
  pdfauthor={Bin Zhang},
  hidelinks,
  pdfcreator={LaTeX via pandoc}}

\title{ECL: A Deterministic Cross-Runtime Execution IR with Replayable
Semantics and Dependency Interface}
\author{Bin Zhang}
\date{2026-06-28}

\begin{document}
\maketitle

\subsection{Abstract}\label{abstract}

Agent runtimes increasingly expose traces, tool calls, callbacks, and
event streams, but these records are usually tied to a specific
framework. This paper presents Execution Compact Layer (ECL), a
deterministic execution intermediate representation for agent systems.
ECL maps runtime traces into four replayable surfaces: \texttt{state},
\texttt{intent}, \texttt{action}, and \texttt{evidence}. The v0.1 system
includes a frozen JSON schema, OpenAI Agents SDK and LangChain trace
adapters, deterministic validation and replay, an embeddable SDK, a
dependency-mode API, and an MCP-shaped local wrapper. Local evaluation
shows stable replay hashes, cross-runtime fixture conversion, schema
validation, and no drift across the locked v0.1 surfaces. ECL is not a
framework, public standard, benchmark, or external adoption claim. It is
a local, reproducible system artifact for representing agent execution
semantics across runtime boundaries.

\subsection{1. Introduction}\label{introduction}

Modern agent systems produce execution traces through framework-specific
surfaces: model inputs, tool calls, callback events, run trees, spans,
and evidence logs. These records are useful inside their host
frameworks, but they are difficult to compare, replay, or cite across
runtimes without rewriting the host runtime or assuming a common
execution model.

ECL addresses this gap by defining a minimal deterministic
representation for agent execution. The goal is not to replace agent
frameworks. The goal is to provide a compact execution IR that can be
generated from heterogeneous traces and replayed into stable local
artifacts.

This paper makes four contributions:

\begin{enumerate}
\def\labelenumi{\arabic{enumi}.}
\tightlist
\item
  A four-surface execution model: \texttt{state}, \texttt{intent},
  \texttt{action}, and \texttt{evidence}.
\item
  A deterministic hash and replay contract for local execution records.
\item
  A cross-runtime mapping surface for OpenAI Agents SDK-style and
  LangChain-style traces.
\item
  A dependency interface and MCP-shaped local wrapper that expose ECL
  without modifying host runtimes.
\end{enumerate}

\subsection{2. Problem Statement}\label{problem-statement}

Agent execution records are fragmented across runtimes. A tool call in
one framework may appear as a span, callback, run node, trace event, or
tool invocation object in another. Existing traces can preserve rich
local detail, but they do not necessarily provide a minimal
cross-runtime representation with deterministic replay semantics.

The problem is:

\begin{Shaded}
\begin{Highlighting}[]
\NormalTok{Given a runtime trace T from framework R,}
\NormalTok{produce a deterministic execution representation E}
\NormalTok{such that E can be validated, replayed, hashed, and cited}
\NormalTok{without modifying R or asserting full{-}fidelity trace preservation.}
\end{Highlighting}
\end{Shaded}

The required properties are:

\begin{itemize}
\tightlist
\item
  Minimality: ECL keeps only the execution surfaces needed for state,
  intent, action, and evidence.
\item
  Determinism: repeated conversion and replay over the same inputs
  produces identical hashes.
\item
  Replayability: a valid ECL object can generate execution trace, replay
  result, and evidence bundle artifacts.
\item
  Loss awareness: adapter mappings can record incomplete preservation
  instead of silently claiming full fidelity.
\item
  Boundary preservation: ECL references host runtime data but does not
  redefine the host runtime.
\end{itemize}

\subsection{3. Related Work}\label{related-work}

\subsubsection{3.1 Runtime Tracing}\label{runtime-tracing}

OpenAI Agents SDK exposes tracing for agent workflows, and LangChain
exposes tracing and callback-style observability for chains, tools, and
model calls. These systems are runtime-specific trace surfaces. ECL
differs by mapping traces into a minimal replayable IR rather than
serving as the primary runtime trace system.

\subsubsection{3.2 Tool and Protocol
Surfaces}\label{tool-and-protocol-surfaces}

The Model Context Protocol (MCP) defines a protocol surface for tools
and context exchange. ECL's MCP-shaped local wrapper is not a conformant
MCP implementation, JSON-RPC transport, published MCP server, or
registry integration. It is a local anchor surface that demonstrates how
an execution representation can be exposed through tool-shaped
functions.

\subsubsection{3.3 Observability and Trace
Standards}\label{observability-and-trace-standards}

OpenTelemetry defines trace data models for distributed systems. ECL is
related in spirit because it treats execution as reviewable trace data,
but ECL is narrower: it targets deterministic agent execution
representation, replay artifacts, and evidence bundles.

\subsubsection{3.4 Intermediate
Representations}\label{intermediate-representations}

Compiler and runtime ecosystems use intermediate representations such as
LLVM IR and WebAssembly to separate source-level systems from execution
targets. ECL uses the same separation principle at a smaller scale: host
runtimes emit or are mapped into a compact execution representation.

\subsubsection{3.5 Canonicalization}\label{canonicalization}

ECL's deterministic hash model depends on sorted-key compact JSON
serialization and SHA-256 hashing. This local serialization is used for
tested hash stability and is not claimed as full RFC 8785/JCS
compliance.

\subsection{4. ECL Model}\label{ecl-model}

An ECL record is a tuple:

\begin{Shaded}
\begin{Highlighting}[]
\NormalTok{E = (S, I, A, V)}
\end{Highlighting}
\end{Shaded}

where:

\begin{itemize}
\tightlist
\item
  \texttt{S} is state: lifecycle, actor, persona, runtime, and
  correlation references.
\item
  \texttt{I} is intent: operation, constraints, expected result, and
  evidence requirements.
\item
  \texttt{A} is action: tool or operation step, execution mode,
  side-effect class, and parameters hash.
\item
  \texttt{V} is evidence: result summary, trace references, event chain,
  policy decisions, and hashes.
\end{itemize}

The machine-readable binding is:

\begin{Shaded}
\begin{Highlighting}[]
\NormalTok{schemas/ecl{-}execution{-}compact{-}layer.schema.json}
\end{Highlighting}
\end{Shaded}

The local semantic source is:

\begin{Shaded}
\begin{Highlighting}[]
\NormalTok{SPEC.md}
\end{Highlighting}
\end{Shaded}

\subsection{5. Formal Execution
Contract}\label{formal-execution-contract}

Let:

\begin{itemize}
\tightlist
\item
  \texttt{T\_r} be a source trace from runtime \texttt{r}.
\item
  \texttt{M\_r} be a deterministic adapter for runtime \texttt{r}.
\item
  \texttt{E} be an ECL object.
\item
  \texttt{L} be a loss report.
\item
  \texttt{P} be the schema and hash validator.
\item
  \texttt{R} be the replay function.
\item
  \texttt{H} be SHA-256 over deterministic sorted-key compact JSON.
\end{itemize}

The adapter contract is:

\begin{Shaded}
\begin{Highlighting}[]
\NormalTok{M\_r(T\_r) {-}\textgreater{} (E, L)}
\end{Highlighting}
\end{Shaded}

The validation contract is:

\begin{Shaded}
\begin{Highlighting}[]
\NormalTok{P(E, schema) {-}\textgreater{} valid | invalid}
\end{Highlighting}
\end{Shaded}

The replay contract is:

\begin{Shaded}
\begin{Highlighting}[]
\NormalTok{R(E) {-}\textgreater{} (execution\_trace.json, evidence\_bundle.json, replay\_result.json)}
\end{Highlighting}
\end{Shaded}

The determinism invariant is:

\begin{Shaded}
\begin{Highlighting}[]
\NormalTok{If T\_r, M\_r, schema, checkpoint\_ref, and generated\_at are unchanged,}
\NormalTok{then H(E) and H(R(E)) remain unchanged.}
\end{Highlighting}
\end{Shaded}

The boundary invariant is:

\begin{Shaded}
\begin{Highlighting}[]
\NormalTok{ECL records execution representation semantics.}
\NormalTok{ECL does not prove that an external side effect occurred.}
\NormalTok{ECL does not make host runtime traces authoritative outside their source systems.}
\end{Highlighting}
\end{Shaded}

\subsection{6. System Design}\label{system-design}

\subsubsection{6.1 Schema Layer}\label{schema-layer}

The schema layer defines the ECL object model and validates required
fields. It is implemented in:

\begin{Shaded}
\begin{Highlighting}[]
\NormalTok{schemas/ecl{-}execution{-}compact{-}layer.schema.json}
\NormalTok{ecl/validator.py}
\end{Highlighting}
\end{Shaded}

\subsubsection{6.2 Adapter Layer}\label{adapter-layer}

The adapter layer maps runtime traces into ECL:

\begin{Shaded}
\begin{Highlighting}[]
\NormalTok{ecl/adapters/openai\_agents.py}
\NormalTok{ecl/adapters/langchain.py}
\NormalTok{external/adoption/ecl\_mapper.py}
\NormalTok{external/adoption/trace\_loader.py}
\end{Highlighting}
\end{Shaded}

The mapping rule is:

{\def\LTcaptype{none} % do not increment counter
\begin{longtable}[]{@{}ll@{}}
\toprule\noalign{}
Runtime source & ECL surface \\
\midrule\noalign{}
\endhead
\bottomrule\noalign{}
\endlastfoot
runtime metadata, actor, correlation & state \\
model input, reasoning, requested operation & intent \\
tool call or run step & action \\
trace events, outputs, result refs & evidence \\
\end{longtable}
}

\subsubsection{6.3 Replay Layer}\label{replay-layer}

The replay layer validates an ECL record and emits deterministic
artifacts:

\begin{Shaded}
\begin{Highlighting}[]
\NormalTok{demo/replay\_demo.py}
\NormalTok{external/adoption/replay\_adapter.py}
\end{Highlighting}
\end{Shaded}

\subsubsection{6.4 SDK and Dependency
Layer}\label{sdk-and-dependency-layer}

The SDK exposes a minimal embedding surface:

\begin{Shaded}
\begin{Highlighting}[]
\NormalTok{sdk/ecl.py}
\NormalTok{sdk/ecl\_dependency.py}
\end{Highlighting}
\end{Shaded}

The dependency API is:

\begin{Shaded}
\begin{Highlighting}[]
\ImportTok{import}\NormalTok{ ecl\_dependency }\ImportTok{as}\NormalTok{ ecl}

\NormalTok{ecl\_object }\OperatorTok{=}\NormalTok{ ecl.wrap(trace)}
\NormalTok{payload }\OperatorTok{=}\NormalTok{ ecl.emit(ecl\_object)}
\NormalTok{result }\OperatorTok{=}\NormalTok{ ecl.verify(ecl\_object)}
\end{Highlighting}
\end{Shaded}

\subsubsection{6.5 MCP-Shaped Local
Wrapper}\label{mcp-shaped-local-wrapper}

The MCP-shaped local wrapper exposes the dependency API through
tool-shaped functions:

\begin{Shaded}
\begin{Highlighting}[]
\NormalTok{mcp/ecl\_tool\_spec.json}
\NormalTok{mcp/ecl\_server\_stub.py}
\end{Highlighting}
\end{Shaded}

This is a local wrapper only. It is not a conformant MCP implementation,
JSON-RPC transport, registry plugin, published MCP server, and not an
external adoption signal.

\subsection{7. Architecture}\label{architecture}

\begin{Shaded}
\begin{Highlighting}[]
\NormalTok{flowchart LR}
\NormalTok{  OA["OpenAI{-}style trace"] {-}{-}\textgreater{} MAP["ECL mapping"]}
\NormalTok{  LC["LangChain{-}style trace"] {-}{-}\textgreater{} MAP}
\NormalTok{  MAP {-}{-}\textgreater{} ECL["ECL object: state / intent / action / evidence"]}
\NormalTok{  ECL {-}{-}\textgreater{} VAL["Validator"]}
\NormalTok{  VAL {-}{-}\textgreater{} REP["Deterministic replay"]}
\NormalTok{  REP {-}{-}\textgreater{} ART["execution\_trace.json / evidence\_bundle.json / replay\_result.json"]}
\NormalTok{  ECL {-}{-}\textgreater{} SDK["SDK dependency API: wrap / emit / verify"]}
\NormalTok{  SDK {-}{-}\textgreater{} MCP["MCP{-}shaped local wrapper"]}
\end{Highlighting}
\end{Shaded}

\subsection{8. Evaluation}\label{evaluation}

\subsubsection{8.1 Method}\label{method}

The evaluation is local and fixture-based. It does not use external APIs
and does not claim third-party validation. The test suite checks:

\begin{itemize}
\tightlist
\item
  ECL schema validation.
\item
  OpenAI-style fixture mapping.
\item
  LangChain-style fixture mapping.
\item
  Replay determinism.
\item
  SDK API stability.
\item
  Dependency API stability.
\item
  MCP-shaped local wrapper safety.
\item
  Citation reproducibility demo stability.
\item
  No drift in locked core files.
\item
  Synthetic cross-runtime trace-corpus conversion, loss reporting, and
  replay stability.
\item
  Field-level mapping coverage over the synthetic trace corpus.
\end{itemize}

\subsubsection{8.2 Current Local Results}\label{current-local-results}

The latest local validation run reported:

\begin{Shaded}
\begin{Highlighting}[]
\NormalTok{python3 {-}m unittest discover {-}s tests}
\NormalTok{Ran 78 tests}
\NormalTok{OK}
\end{Highlighting}
\end{Shaded}

The synthetic trace-corpus evaluation reported:

\begin{Shaded}
\begin{Highlighting}[]
\NormalTok{python3 experiments/evaluate\_trace\_corpus.py}
\NormalTok{case\_count=12}
\NormalTok{by\_runtime=\{"langchain": 6, "openai": 6\}}
\NormalTok{valid\_count=12}
\NormalTok{deterministic\_count=12}
\NormalTok{loss\_expectation\_met\_count=12}
\NormalTok{loss\_detected\_count=3}
\NormalTok{surface\_coverage=\{"action": 12, "evidence": 12, "intent": 12, "state": 12\}}
\NormalTok{surface\_coverage\_by\_runtime=\{"langchain": \{"action": 6, "evidence": 6, "intent": 6, "state": 6\}, "openai": \{"action": 6, "evidence": 6, "intent": 6, "state": 6\}\}}
\NormalTok{loss\_type\_counts=\{"semantic": 9, "structural": 3\}}
\NormalTok{evaluation\_hash=sha256:2434da21056811e7eacf1ffac6944d5420ddeb2aa783f74423326a2b54a33974}
\end{Highlighting}
\end{Shaded}

The corpus contains local synthetic stress cases for complete traces,
missing fields, unknown fields, failure status, timestamp drift, nested
runs, and implicit tool events. It is not production trace evidence,
third-party validation, or a benchmark result.

The negative validation-matrix evaluation reported:

\begin{Shaded}
\begin{Highlighting}[]
\NormalTok{python3 experiments/evaluate\_validation\_matrix.py}
\NormalTok{baseline\_valid=true}
\NormalTok{case\_count=8}
\NormalTok{expected\_invalid\_count=8}
\NormalTok{invalid\_detected\_count=8}
\NormalTok{expectation\_met\_count=8}
\NormalTok{evaluation\_hash=sha256:00fe0ee8952988f7460280b40554889a890c93f9d181a6c9d8ee8b0194b031c0}
\end{Highlighting}
\end{Shaded}

The validation matrix mutates a valid ECL record to test rejection
behavior for missing required surfaces, wrong schema version, tampered
canonical hash, source-hash mismatch, empty event chains, additional
root properties, invalid action modes, and missing evidence hashes.

The field-level mapping coverage evaluation reported:

\begin{Shaded}
\begin{Highlighting}[]
\NormalTok{python3 experiments/evaluate\_mapping\_coverage.py}
\NormalTok{case\_count=12}
\NormalTok{total\_source\_fields=81}
\NormalTok{direct\_mapped\_field\_count=80}
\NormalTok{source\_hash\_only\_field\_count=1}
\NormalTok{loss\_missing\_field\_count=4}
\NormalTok{evaluation\_hash=sha256:8a8b820ecbdd1b4e88a2d8e07b05be2d479add0f0c6c3265292cc5a86763e43a}
\end{Highlighting}
\end{Shaded}

This mapping coverage evaluation records whether each top-level source
trace field is projected into an ECL surface or retained only through
the source trace hash. It is synthetic corpus evidence, not a benchmark
or external adoption signal.

The citation reproducibility demo reported:

\begin{Shaded}
\begin{Highlighting}[]
\NormalTok{python3 examples/citation\_repro\_demo.py}
\NormalTok{all\_valid=true}
\NormalTok{all\_deterministic=true}
\NormalTok{result\_hash=sha256:8684fa8963ff9ad55c36eed6b89eb15992d1b05566aa1c79ced754fac67e1cdd}
\end{Highlighting}
\end{Shaded}

The MCP-shaped local wrapper self-check reported valid deterministic
execution with:

\begin{Shaded}
\begin{Highlighting}[]
\NormalTok{verification\_hash=sha256:bdcf2ddb30a0e5975782f77fa86415226c46652bbd5490a667a2f2106dd65a5f}
\end{Highlighting}
\end{Shaded}

\subsubsection{8.3 Locked Core Surfaces}\label{locked-core-surfaces}

The current local evaluation keeps these core surfaces unchanged:

\begin{Shaded}
\begin{Highlighting}[]
\NormalTok{schemas/ecl{-}execution{-}compact{-}layer.schema.json}
\NormalTok{sdk/ecl.py}
\NormalTok{sdk/ecl\_dependency.py}
\NormalTok{mcp/ecl\_tool\_spec.json}
\NormalTok{mcp/ecl\_server\_stub.py}
\end{Highlighting}
\end{Shaded}

\subsubsection{8.4 Interpretation}\label{interpretation}

The results support a narrow claim: ECL v0.1 is locally reproducible
over the included fixtures, synthetic trace corpus, and deterministic
entrypoints. The results do not support claims of ecosystem adoption,
production reliability, third-party validation, or benchmark
superiority.

\subsection{9. Discussion}\label{discussion}

ECL is most useful when execution records must be compared, replayed, or
cited across runtime boundaries. The main design choice is to preserve a
small execution surface rather than reproduce every runtime-specific
field. This makes the representation compact but requires explicit loss
reports when mappings are incomplete.

ECL also separates representation from execution. A host runtime
executes. ECL records and replays the execution representation. This
separation keeps the dependency interface non-invasive and avoids
modifying host frameworks.

The MCP-shaped local wrapper is intentionally conservative. It
demonstrates tool-shaped surface readability but does not claim MCP
protocol conformance or adoption.

\subsection{10. Limitations}\label{limitations}

\begin{itemize}
\tightlist
\item
  The evaluation uses local fixtures, not production traces.
\item
  The system does not prove full-fidelity trace preservation.
\item
  The current adapters target OpenAI-style and LangChain-style traces
  only.
\item
  The MCP surface is a local stub, not a published MCP server.
\item
  There is no external adoption signal in this version.
\item
  There is no benchmark suite or leaderboard result.
\end{itemize}

\subsection{11. Conclusion}\label{conclusion}

ECL v0.1 demonstrates that agent execution traces can be mapped into a
deterministic, replayable, hash-stable execution IR with a minimal
dependency interface. The system is locally validated and
citation-ready, but the paper-level claim remains deliberately narrow:
ECL is a reproducible execution representation stack, not a framework,
formal standard, production deployment, third-party validation result,
or external adoption result.

\subsection{References}\label{references}

\begin{itemize}
\tightlist
\item
  OpenAI Agents SDK tracing documentation:
  https://openai.github.io/openai-agents-python/tracing/
\item
  LangChain tracing documentation:
  https://docs.langchain.com/langsmith/tracing
\item
  Model Context Protocol specification:
  https://modelcontextprotocol.io/specification/
\item
  OpenTelemetry trace specification:
  https://opentelemetry.io/docs/specs/otel/trace/
\item
  RFC 8785 JSON Canonicalization Scheme:
  https://www.rfc-editor.org/rfc/rfc8785
\item
  LLVM Language Reference Manual: https://llvm.org/docs/LangRef.html
\item
  WebAssembly Core Specification:
  https://webassembly.github.io/spec/core/
\end{itemize}

\end{document}
